\paragraph{\trAdditive{}扩大。}\label{sub:additive-enlargement}
构造 $Tw\scr{A}$ 的第一步，是\textbf{\trAdditive{}扩大}（additive enlargement）$\Sigma\scr{A}$。
设 $\scr{A}$ 是一个\trNonunital{} \trACategory。$\Sigma\scr{A}$ 的\trObject{}是三元组
$X=(I,\{X^i\},\{V^i\})$，写成更直观的形式\trDirectSum{}即
\begin{equation}\label{eq:sigma-formal-direct-sum}
X=\bigoplus_{i\in I}V^i\otimes X^i,
\end{equation}
其中 $I$ 是一个有限集；$\{X^i\}_{i\in I}$ 是 $\scr{A}$ 的一族\trObject，$\{V^i\}_{i\in I}$
是与之相应的一族有限维\trGraded\trVectorSpace，称为\textbf{\trMultiplicity{}空间}（multiplicity spaces）。
两个这种\trDirectSum{}之间的\trMorphism，由各\trDirectSummand{}之间的\trMorphism{}与相应\trVectorSpace{}上的%
\trLinear{}映射作适当\trTensorProduct{}拼成：
\begin{equation}\label{eq:sigma-morphism-spaces}
\hom_{\Sigma\scr{A}}\Bigl(\bigoplus_{i\in I_0}V_0^i\otimes X_0^i,
\bigoplus_{j\in I_1}V_1^j\otimes X_1^j\Bigr)
=\bigoplus_{ij}\hom_{\KK}(V_0^i,V_1^j)\otimes\hom_{\scr{A}}(X_0^i,X_1^j),
\end{equation}
配以自然的\trGrading。我们把\trMorphism{}写成矩阵 $a=(a^{ji})$，元素分别以 $j\in I_1$ 与
$i\in I_0$ 标号，每个元素形如
$a^{ji}=\sum_k\phi^{jik}\otimes x^{jik}\in\hom_{\KK}(V_0^i,V_1^j)
\otimes\hom_{\scr{A}}(X_0^i,X_1^j)$。在书写公式时，我们常限于该和只含一项的情形
$a^{ji}=\phi^{ji}\otimes x^{ji}$，其\trLinear{}性则默认成立。有了这些准备，就可以定义
$\Sigma\scr{A}$ 上的\trComposition{}映射了。它们把 $\scr{A}$ 的\trComposition{}与%
\trLinear{}映射的通常\trComposition{}结合起来，正如\cref{thm:homotopy-rational-interpretation}
中的\trTensorProduct{}那样：
\begin{equation}\label{eq:sigma-composition}
\begin{split}
\mu^d_{\Sigma\scr{A}}(a_d,\ldots,a_1)^{i_d,i_0}
&=\sum_{i_1,\ldots,i_{d-1}}(-1)^{\lhd}\,
\phi_d^{i_d,i_{d-1}}\cdots\phi_1^{i_1,i_0}\otimes\\
&\quad\otimes\mu^d_{\scr{A}}\bigl(x_d^{i_d,i_{d-1}},\ldots,x_1^{i_1,i_0}\bigr),
\end{split}
\end{equation}
其中
$\lhd=\sum_{p<q}\bigl|\phi_p^{i_p,i_{p-1}}\bigr|\cdot
\bigl(\bigl|x_q^{i_q,i_{q-1}}\bigr|-1\bigr)$，
而 $a_k=(a_k^{ji})=(\phi_k^{ji}\otimes x_k^{ji})$。我们将一劳永逸地取定一个单元素集
$\{*\}$，并把 $\scr{A}$ 的\trObject{}等同于三元组
$(\{*\},\{X^*=X\},\{V^*=\KK\})$，其中 $\KK$ 被置于零次。这样便把 $\scr{A}$ 作为一个%
\trFull{} \trAinfty-子\trCategory{}\trEmbedding{}到 $\Sigma\scr{A}$。
